\documentclass[10pt,a4paper]{amsart}
\usepackage{amsmath,amssymb,amsthm,graphicx}
\usepackage{verbatim}
\newenvironment{reference}{\comment}{\endcomment}
%\newenvironment{reference}{}{}
\newenvironment{slogan}{\comment}{\endcomment}
\newenvironment{history}{\comment}{\endcomment}

% For commutative diagrams we use Xy-pic
\usepackage[all]{xy}

% We use 2cell for 2-commutative diagrams.
\xyoption{2cell}
% Theorem environments.
%
\theoremstyle{plain}
\newtheorem{theorem}[subsection]{Theorem}
\newtheorem{proposition}[subsection]{Proposition}
\newtheorem{lemma}[subsection]{Lemma}

\theoremstyle{definition}
\newtheorem{definition}[subsection]{Definition}
\newtheorem{example}[subsection]{Example}
\newtheorem{exercise}[subsection]{Exercise}
\newtheorem{situation}[subsection]{Situation}

\theoremstyle{remark}
\newtheorem{remark}[subsection]{Remark}
\newtheorem{remarks}[subsection]{Remarks}

\numberwithin{equation}{subsection}

% Macros
%
\def\lim{\mathop{\mathrm{lim}}\nolimits}
\def\colim{\mathop{\mathrm{colim}}\nolimits}
\def\Spec{\mathop{\mathrm{Spec}}}
\def\Hom{\mathop{\mathrm{Hom}}\nolimits}
\def\Ext{\mathop{\mathrm{Ext}}\nolimits}
\def\SheafHom{\mathop{\mathcal{H}\!\mathit{om}}\nolimits}
\def\SheafExt{\mathop{\mathcal{E}\!\mathit{xt}}\nolimits}
\def\Sch{\mathit{Sch}}
\def\Mor{\mathop{\mathrm{Mor}}\nolimits}
\def\Ob{\mathop{\mathrm{Ob}}\nolimits}
\def\Sh{\mathop{\mathit{Sh}}\nolimits}
\def\NL{\mathop{N\!L}\nolimits}
\def\CH{\mathop{\mathrm{CH}}\nolimits}
\def\proetale{{pro\text{-}\acute{e}tale}}
\def\etale{{\acute{e}tale}}
\def\QCoh{\mathit{QCoh}}
\def\Ker{\mathop{\mathrm{Ker}}}
\def\Im{\mathop{\mathrm{Im}}}
\def\Coker{\mathop{\mathrm{Coker}}}
\def\Coim{\mathop{\mathrm{Coim}}}

% Boxtimes
%
\DeclareMathSymbol{\boxtimes}{\mathbin}{AMSa}{"02}

%
% Macros for moduli stacks/spaces
%
\def\QCohstack{\mathcal{QC}\!\mathit{oh}}
\def\Cohstack{\mathcal{C}\!\mathit{oh}}
\def\Spacesstack{\mathcal{S}\!\mathit{paces}}
\def\Quotfunctor{\mathrm{Quot}}
\def\Hilbfunctor{\mathrm{Hilb}}
\def\Curvesstack{\mathcal{C}\!\mathit{urves}}
\def\Polarizedstack{\mathcal{P}\!\mathit{olarized}}
\def\Complexesstack{\mathcal{C}\!\mathit{omplexes}}
% \Pic is the operator that assigns to X its picard group, usage \Pic(X)
% \Picardstack_{X/B} denotes the Picard stack of X over B
% \Picardfunctor_{X/B} denotes the Picard functor of X over B
\def\Pic{\mathop{\mathrm{Pic}}\nolimits}
\def\Picardstack{\mathcal{P}\!\mathit{ic}}
\def\Picardfunctor{\mathrm{Pic}}
\def\Deformationcategory{\mathcal{D}\!\mathit{ef}}
\usepackage{cleveref}
\crefname{originallemma}{original lemma}{original lemmas}
\Crefname{originallemma}{Original lemma}{Original lemmas}
\makeatletter
\newlabel{lemma-NAK}{{00DV}{source}{}{sourcepointer-1}{}}
\newlabel{lemma-NAK@cref}{{[originallemma][0][]00DV}{[1][0][]source}{}{}{}}
\makeatother
\title{A regular sequence over a local Noetherian ring stays regular after any permutation}
\author{The Stacks Project Authors}
\date{stacks-project:algebra.tex:local-snapshot-20261012; GNU FDL 1.2 or later}
\begin{document}
\maketitle
\noindent\textit{Editorial note.} All delivered mathematical text is the author's own, copied from the pinned Stacks Project chapter file algebra.tex; no statement or proof step was rewritten, reordered or removed.\par\medskip
\noindent\textit{Editorial note.} Preamble provenance, as required for this source: the frozen snapshot stages the chapter's own preamble file beside it (preamble.tex, 217 lines, 8116 bytes, the same shared Stacks-wide preamble the published chow and etale-cohomology cards use); the first line of the pinned chapter file is \textbackslash{}input\{preamble\}, so no preamble text lives inside the chapter itself. The card takes that file's lines 8-18 (the verbatim package and the three \textbackslash{}comment-based environments reference/slogan/history at lines 9-12, section 13.2, together with the Xy-pic setup) and its lines 154 to the runtime end of the file (the native theorem declarations and mathematical macros). Not one character was changed, no package was added, no macro was defined, and the excerpt serves only as preamble, never as a span coordinate.\par\medskip
\noindent\textit{Editorial note.} The same-chapter pointers are located at their \textbackslash{}label lines inside the pinned chapter file; the cross-chapter pointers are located in the project's own tag list staged in the snapshot (sources/stacks-chow/tags.txt), at the verbatim lines carrying their canonical tags.\par\medskip
\noindent\textit{Editorial note.} The counted proof refers to Nakayama's lemma by label; it is recorded as a navigation stub pointed at its \textbackslash{}label line, and the argument consuming it stays inside the counted span.\par\medskip
\noindent\textit{Editorial note.} The counted spans contain no \textbackslash{}cite, so the card carries no bibliography.\par\medskip
\noindent\textit{Allowed standard prerequisites (editorial).}\begin{itemize}
\item Nakayama's lemma. Invoked in the counted proof at line 16872. Turns surjectivity of multiplication by x\_1 on the kernel K into K = 0.
\item Regularity on a quotient by the previous elements. Used at the induction step of the counted proof (lines 16887-16889). Reduces a transposition in the middle of the sequence to the two-term case over a quotient module.
\end{itemize}\medskip
\ifdefined\hypertarget\hypertarget{sourcepointer-1}{}\fi
\noindent\textit{Source reference.} Original lemma 00DV: sources/stacks-algebra/pinned-original-source.tex, source lines 3833--3833. Navigation pointer for \textbackslash{}ref\{lemma-NAK\} used inside the counted span (line 16872); the argument that consumes it is retained in the span.\par\medskip
\begin{lemma}
\label{lemma-permute-xi}
Let $R$ be a local Noetherian ring.
Let $M$ be a finite $R$-module.
Let $x_1, \ldots, x_c$ be an $M$-regular sequence.
Then any permutation of the $x_i$ is a regular
sequence as well.
\end{lemma}

\begin{proof}
First we do the case $c = 2$.
Consider $K \subset M$ the kernel of $x_2 : M \to M$.
For any $z \in K$ we know that $z = x_1 z'$
for some $z' \in M$ because
$x_2$ is a nonzerodivisor on $M/x_1M$.
Because $x_1$ is a nonzerodivisor on $M$ we see that $x_2 z' = 0$
as well. Hence $x_1 : K \to K$ is surjective.
Thus $K = 0$ by Nakayama's Lemma \ref{lemma-NAK}.
Next, consider multiplication by $x_1$ on $M/x_2M$.
If $z \in M$ maps to an element $\overline{z} \in M/x_2M$
in the kernel of this map, then $x_1 z = x_2 y$ for some $y \in M$.
But then since $x_1, x_2$ is a regular sequence we see that
$y = x_1 y'$ for some $y' \in M$. Hence $x_1 ( z - x_2 y' ) =0$
and hence $z = x_2 y'$ and hence $\overline{z} = 0$ as desired.

\medskip\noindent
For the general case, observe that any permutation is
a composition of transpositions of adjacent indices.
Hence it suffices to prove that
$$
x_1, \ldots, x_{i-2}, x_i, x_{i-1}, x_{i + 1}, \ldots, x_c
$$
is an $M$-regular sequence. This follows from the case we
just did applied to the module $M/(x_1, \ldots, x_{i-2})$
and the length $2$ regular sequence $x_{i-1}, x_i$.
\end{proof}

\end{document}
